\begin{table}[htbp]
\centering
\small
\caption{Viés de variantes de prompt (openai/gpt-4o-mini): fidelidade, direção de $\hat{W}$ e consistência.}
\label{tab:bloco1_vies_variantes_prompt_gpt4mini}
\begin{tabular}{llrrrrrrr}
\toprule
Variante & Regime & Fid.\ A & Acur.\ LLM vs real A & $w_1/w_2$ & Consist.\ B & $\kappa$ B & Consist.\ C & $\kappa$ C \\
\midrule
chain-of-thought & zero-shot & 71,1 (67,3–75,3) & 36,9 & 0,00 & 62,3 & 0,157 & 75,7 & 0,513 \\
\textbf{padrão} & \textbf{zero-shot} & \textbf{81,3 (78,0–85,3)} & \textbf{32,4} & \textbf{0,09} & \textbf{70,7} & \textbf{0,348} & \textbf{88,0} & \textbf{0,745} \\
geométrica & zero-shot & 78,0 (72,0–82,0) & 48,4 & 0,48 & 64,7 & 0,112 & 76,7 & 0,413 \\
tabular & zero-shot & 65,8 (56,0–74,0) & 47,6 & 0,10 & 63,3 & 0,150 & 72,3 & 0,186 \\
\midrule
chain-of-thought & 10-shot & 71,1 (67,3–75,3) & 36,9 & 0,00 & 93,7 & 0,873 & 88,5 & 0,767 \\
\textbf{padrão} & \textbf{10-shot} & \textbf{81,3 (78,0–85,3)} & \textbf{32,4} & \textbf{0,09} & \textbf{82,2} & \textbf{0,624} & \textbf{91,9} & \textbf{0,834} \\
geométrica & 10-shot & 78,0 (72,0–82,0) & 48,4 & 0,48 & 85,6 & 0,707 & 93,0 & 0,859 \\
tabular & 10-shot & 65,8 (56,0–74,0) & 47,6 & 0,10 & 95,9 & 0,919 & 85,6 & 0,684 \\
\bottomrule
\multicolumn{9}{p{0.95\linewidth}}{\footnotesize Problema A e transferência para B e C por variantes de prompt; os demais fatores ficam na configuração base. Médias das sementes (fidelidade com mín–máx). Fid.\ A = fidelidade (vs LLM) no A; Acur.\ LLM vs real A = acurácia (vs rótulo real) do LLM no A, contexto; $w_1/w_2$ = razão da métrica estimada (Perceptron). Fid.\ A, Acur.\ LLM vs real A e $w_1/w_2$ vêm da Fase A, sempre zero-shot (iguais nas linhas zero e 10-shot). $\kappa$ = Cohen's $\kappa$ (faixas de Landis e Koch, 1977). Em negrito: a configuração de referência. Em zero-shot (e com nomes de classe neutros) a correspondência entre nome e classe é arbitrária: a acurácia (vs rótulo real) do LLM pode sair abaixo de 50\%; o valor invariante à permutação é max($a$, 1 $-$ $a$). 10-shot: exemplos rotulados pela métrica zero-shot do A (H2). Os exemplos few-shot de B e C são rotulados por $\hat{W}$ do A com os centróides do A; a consistência usa os centróides locais, recalculados pelos rótulos do LLM no próprio problema.} \\
\end{tabular}
\end{table}
